\section{PART I --- FOUNDATIONS}

\subsection{1. Errors, Conditioning, and Floating Point}
\begin{itemize}
\item IEEE arithmetic; machine epsilon; rounding models
\item Condition number of a problem (not just a matrix)
\item Forward and backward error analysis; Wilkinson's program
\item Catastrophic cancellation; examples of unstable algorithms
\end{itemize}

\subsection{2. Polynomial Approximation and Interpolation}
\begin{itemize}
\item Lagrange and Newton interpolation; divided differences
\item Chebyshev nodes; Runge's phenomenon; optimal point selection
\item Spline interpolation: cubic splines, B-splines, tension splines
\item Barycentric interpolation formula; numerical stability
\end{itemize}

\subsection{3. Approximation Theory}
\begin{itemize}
\item Best approximation in normed spaces; existence, uniqueness
\item Chebyshev equioscillation theorem; Remez algorithm
\item Orthogonal polynomials: Legendre, Chebyshev, Hermite, Laguerre
\item Rational approximation; Pad\'e approximants
\end{itemize}

\subsection{4. Numerical Differentiation and Integration}
\begin{itemize}
\item Finite differences; Richardson extrapolation; accuracy
\item Newton-Cotes formulas; error analysis
\item Gaussian quadrature: nodes, weights, optimality
\item Adaptive quadrature; Gauss-Kronrod; oscillatory integrals
\end{itemize}

\section{PART II --- CORE ALGORITHMS}

\subsection{5. Direct Methods for Linear Systems}
\begin{itemize}
\item LU factorization; pivoting; stability analysis
\item Cholesky for SPD; LDLT decomposition; tridiagonal systems
\item QR factorization: Householder, Givens, modified Gram-Schmidt
\item Least squares; pseudoinverse; rank-deficient problems; regularization
\end{itemize}

\subsection{6. Iterative Methods for Linear Systems}
\begin{itemize}
\item CG (SPD), MINRES (symmetric), GMRES (general), BiCGSTAB
\item Convergence theory; spectral analysis; polynomial approximation view
\item Preconditioning: ILU, incomplete Cholesky, algebraic multigrid (conceptual)
\item Choosing a solver; practical considerations
\end{itemize}

\subsection{7. Eigenvalue Algorithms}
\begin{itemize}
\item Power method; inverse iteration; shift-and-invert
\item QR algorithm with shifts and deflation; convergence theory
\item Symmetric QR with Wilkinson shift; divide-and-conquer
\item Krylov methods: Lanczos (symmetric), Arnoldi (general); thick restart
\end{itemize}

\subsection{8. Nonlinear Equations and Fixed Points}
\begin{itemize}
\item Bisection; secant method; Newton's method; convergence analysis
\item Newton's method in $\mathbb{R}^n$; convergence, local vs.\ global
\item Fixed-point iterations; Banach contraction; Aitken-Steffensen acceleration
\item Homotopy and path-following methods; applications
\end{itemize}

\subsection{9. Numerical ODEs}
\begin{itemize}
\item One-step methods: Euler, Runge-Kutta, order conditions
\item Multi-step methods: Adams, BDF; zero-stability; Dahlquist barrier
\item Stiff systems; implicit methods; A-stability; L-stability
\item Event detection; discontinuities; delay equations
\end{itemize}

\subsection{10. Numerical PDEs}
\begin{itemize}
\item Finite differences: elliptic, parabolic, hyperbolic; stability analysis
\item Finite element methods: weak form, assembly, error analysis
\item Spectral methods: Fourier, Chebyshev, pseudospectral collocation
\item Multigrid: V-cycle, smoothers, convergence theory
\end{itemize}

\section{PART III --- MODERN METHODS}

\subsection{11. Optimization Algorithms}
\begin{itemize}
\item Gradient descent and line search; Newton and quasi-Newton (BFGS, L-BFGS)
\item Interior point methods; barrier approach
\item Nonlinear least squares: Gauss-Newton, Levenberg-Marquardt
\item Derivative-free optimization: Nelder-Mead, pattern search (brief)
\end{itemize}

\subsection{12. Automatic Differentiation}
\begin{itemize}
\item Forward mode via dual numbers; mathematical formalization
\item Reverse mode (backpropagation) via computational graphs; chain rule at scale
\item Jacobian-vector products (JVPs) and vector-Jacobian products (VJPs)
\item Higher-order AD: forward-over-reverse; Hessian-vector products
\item JAX, PyTorch autograd, Zygote (Julia) as implementations
\item Adjoint method for PDE-constrained optimization; connection to Course 3 §22
\end{itemize}

\subsection{13. Monte Carlo and Quasi-Monte Carlo Methods}
\begin{itemize}
\item Classical Monte Carlo; variance estimation; confidence intervals
\item Variance reduction: importance sampling, control variates, stratification
\item Markov chain Monte Carlo: Metropolis-Hastings, Gibbs sampler
\item Hamiltonian Monte Carlo; MALA
\item Quasi-Monte Carlo: Halton and Sobol sequences; Koksma-Hlawka inequality
\item Applications: Bayesian computation, uncertainty quantification, rare events
\end{itemize}

\subsection{14. High-Dimensional Numerical Methods}
\begin{itemize}
\item Curse of dimensionality: what breaks and what survives
\item Sparse grids: Smolyak construction; dimension-adaptive variants
\item Tensor decompositions in numerical PDE: TT/MPS, Tucker
\item Sparse polynomial chaos; stochastic collocation for UQ
\item Importance for quantum chemistry, financial mathematics, statistical learning
\end{itemize}

\subsection{15. Randomized Linear Algebra}
\begin{itemize}
\item Randomized SVD (Halko-Martinsson-Tropp); single-pass algorithms
\item Johnson-Lindenstrauss lemma; random projections; sketching
\item Randomized least squares; sketch preconditioners
\item Nystr\"om approximation for kernel matrices; applications
\end{itemize}

\subsection{16. Structured Matrices and Fast Algorithms}
\begin{itemize}
\item Toeplitz and circulant matrices; FFT as diagonalization
\item Hierarchical matrices (H-matrices): structure, arithmetic
\item Fast multipole method (conceptual)
\item Applications: structured covariance, integral equations
\end{itemize}

\subsection{17. Matrix Functions}
\begin{itemize}
\item Matrix exponential: scaling-squaring with Pad\'e; Krylov methods for large systems
\item Other functions: square root, logarithm, sign, cosine
\item Fr\'echet derivative of matrix functions; algorithmic differentiation
\item Applications: network centrality, Markov chains, differential equations
\end{itemize}

\subsection{18. Probabilistic Numerics}
\begin{itemize}
\item Bayesian quadrature: integration with uncertainty quantification
\item Gaussian process models for ODE solutions; probabilistic linear solvers
\item Connections between numerical algorithms and Bayesian inference
\item Uncertainty-aware scientific computing
\end{itemize}
